% Drawing 74/03-4 -- folder 74, batch 03, archive page 44 (PDF page 45).
% Transcription: à gauche, deux hexagones de sommets $x_1,\dots,x_6$, avec les paires $\{1,2\}$, $\{2,5\}$, $\{1,6\}$, $\{1,5\}$, $\{3,4\}$, $\{2,4\}$, $\{2,6\}$, $\{1,5\}$, $\{1,6\}$, $\{2,4\}$ écrites à côté.
% Kind: two hexagons with chords, vertex labels xi_ij and x_i. Method 2 (schematic redraw from the facsimile). Status: partial.
% Notes: Upper hexagon: vertices labelled xi12 (top left), xi25 (top), an overwritten xi at the top right (not given), xi16 (right), xi24 (bottom), xi3. (left, the second index not legible; xi34 by the transcription's list of pairs); sides as on the page, with the chords left - top right and bottom - top right (each doubled on the page by a second parallel stroke, drawn once). Lower hexagon: vertices x2? (top left), x4 (top), x3 (right), x5 (bottom right), x6 (bottom left), x1 (left), each followed on the page by « = » and an overwritten xi, left out; sides, the chords left - top and top left - bottom left, and an inner point joined to the two bottom vertices (bold on the page). To the right, the pairs xi14, xi16, xi25, xi16, xi24 written in a column, placed as text. Every edge is drawn between the same vertex points. Still differs: the vertex labels of the lower hexagon and the overwritten top right one of the upper hexagon are readings; the doubled strokes are drawn once; the short strokes above the lower hexagon (« 1 »-like marks) are not drawn. Disagreement: the transcription reads the labels as pairs {1,2}, {2,5}, ... written beside two hexagons on x1..x6; on the page the upper hexagon's vertices themselves carry xi12, xi25, ..., and the lower one carries x1..x6. Redrawn 2026-11-09 (issue 128): to check again.
% Coordinates: 1 unit = 1 mm = 1/100 of the crop width (bbox in page fractions: [0.08, 0.06, 0.4, 0.3]).
% Redrawn by Claude (claude-opus-5-5), issue #73, 2026-11-09; not checked by a human; not compiled (no TeX in the session).
% Rights: reproduces part of the work -- behind RENDER_TRANSCRIPTS, never published by itself.
\documentclass[tikz,border=3pt]{standalone}
\usepackage{amsmath,amssymb}
\usetikzlibrary{patterns,arrows.meta}
\tikzset{pencil/.style={black!45}, struck/.style={black!35}, mark/.style={circle,draw,inner sep=1.3pt}, pt/.style={circle,fill,inner sep=0.9pt}, opt/.style={circle,draw,fill=white,inner sep=1.1pt}}
\begin{document}
\begin{tikzpicture}[x=1mm,y=1mm,line cap=round,line join=round,font=\scriptsize]
\draw (28.5,-18.6) -- (41.9,-18.1);
\draw (41.9,-18.1) -- (59.3,-18.1);
\draw (59.3,-18.1) -- (58.2,-29.3);
\draw (58.2,-29.3) -- (31.3,-47.2);
\draw (31.3,-47.2) -- (19,-36);
\draw (19,-36) -- (28.5,-18.6);
\draw (19,-36) -- (59.3,-18.1);
\draw (31.3,-47.2) -- (59.3,-18.1);
\node[inner sep=1pt] at (24,-15.3) {\small $\xi_{12}$};
\node[inner sep=1pt] at (45.3,-14.7) {\small $\xi_{25}$};
\node[inner sep=1pt] at (62.6,-32.6) {\small $\xi_{16}$};
\node[inner sep=1pt] at (32.4,-52.2) {\small $\xi_{24}$};
\node[inner sep=1pt] at (14.5,-37.1) {\small $\xi_{3\cdot}$};
\draw (20.6,-64.2) -- (37.2,-63.5);
\draw (37.2,-63.5) -- (46.8,-80.2);
\draw (46.8,-80.2) -- (37.2,-96.8);
\draw (37.2,-96.8) -- (21.8,-94.2);
\draw (21.8,-94.2) -- (11,-78.9);
\draw (11,-78.9) -- (20.6,-64.2);
\draw (11,-78.9) -- (37.2,-63.5);
\draw (20.6,-64.2) -- (21.8,-94.2);
\draw (31.4,-85.3) -- (21.8,-94.2);
\draw (31.4,-85.3) -- (37.2,-96.8);
\draw[very thick] (31.4,-85.3) -- (37.2,-96.8);
\node[inner sep=1pt] at (16.1,-61) {\small $x_2$};
\node[inner sep=1pt] at (40.4,-59.7) {\small $x_4$};
\node[inner sep=1pt] at (52.6,-80.8) {\small $x_3$};
\node[inner sep=1pt] at (41.7,-101.3) {\small $x_5$};
\node[inner sep=1pt] at (19.3,-100.6) {\small $x_6$};
\node[inner sep=1pt] at (6.5,-77.6) {\small $x_1$};
\node[inner sep=1pt] at (88,-37) {\small $\xi_{14}$};
\node[inner sep=1pt] at (66,-44) {\small $\xi_{25}$};
\node[inner sep=1pt] at (66,-62) {\small $\xi_{16}$};
\node[inner sep=1pt] at (86,-83) {\small $\xi_{24}$};
\end{tikzpicture}
\end{document}
